\part{Appendices}

\chapter{LaTeX notation reference}
\label{app:notation}
This appendix summarises notation used in equations and implementation checks.


{\small
\begin{longtable}{L{0.22\textwidth}L{0.68\textwidth}}
\caption{Core notation.}\\
\toprule
Symbol & Meaning \\
\midrule
\endfirsthead
\toprule
Symbol & Meaning \\
\midrule
\endhead
\(B,T,S,L\) & Batch size, current input length, visible cache length, and number of layers. \\
\(V,P,H,F\) & Vocabulary size, maximum position length, hidden size, and feed-forward width. \\
\(H_q,H_{kv},d_h,g\) & Query heads, key/value heads, per-head dimension, and GQA group size \(g=H_q/H_{kv}\). \\
\(I\) & Token-id tensor in \(\N^{B\times T}\). \\
\(X^{(\ell)}\) & Hidden states entering block \(\ell\), with shape \([B,T,H]\). \\
\(K_\ell,V_\ell\) & Layer-\(\ell\) KV-cache tensors with logical shape \([B,H_{kv},S,d_h]\). \\
\(Z\) & Logits with shape \([B,T,V]\). \\
\bottomrule
\end{longtable}}


\begin{intuitionbox}{Intuition}
Notation reduces ambiguity.
\end{intuitionbox}

\begin{definitionbox}{Mathematical or contract statement}
The same symbols are used in formulas, tensor shape tables, and validators.
\end{definitionbox}

\begin{implementationbox}{Implementation interpretation}
Update notation before adding a new module.
\end{implementationbox}

\begin{verificationbox}{How to verify this works}
\begin{itemize}
\item Check shapes and dtypes at the boundary.
\item Compare deterministic outputs against PyTorch golden vectors.
\item Record tolerance, seed, model artifact, and backend.
\item Do not benchmark until validation passes.
\end{itemize}
\end{verificationbox}

\begin{warningbox}{Common failure modes}
\begin{itemize}
\item Letting framework defaults choose dtype or layout.
\item Testing only a batch size of one.
\item Depending on upstream checkpoint names in runtime code.
\item Treating benchmark output as validation.
\end{itemize}
\end{warningbox}

\begin{chapterexercises}
\item State the central invariant in one sentence.\\
\textit{Hint.} Use the conceptual guide vocabulary.
\item Construct a toy example with small shapes.\\
\textit{Hint.} Use $B=2$, $T=3$, $H=8$ where possible.
\item Name a validation test that would catch a plausible bug.\\
\textit{Hint.} Prefer generated golden vectors to handwritten long arrays.
\end{chapterexercises}

\chapter{Tensor shape cheat sheet}
\label{app:shape-cheat}
This appendix collects common tensor shapes.


{\small
\begin{longtable}{L{0.28\textwidth}L{0.28\textwidth}L{0.34\textwidth}}
\caption{Tensor shape cheat sheet.}\\
\toprule
Tensor & Shape & Notes \\
\midrule
\endfirsthead
\toprule
Tensor & Shape & Notes \\
\midrule
\endhead
Input ids & $[B,T]$ & Integer token ids. \\
Hidden states & $[B,T,H]$ & Residual stream. \\
Queries & $[B,H_q,T,d_h]$ & Per-head query tensor. \\
Keys/values & $[B,H_{kv},S,d_h]$ & Cached attention state by layer. \\
Attention output by head & $[B,H_q,T,d_h]$ & Concatenate heads back to $[B,T,H]$. \\
Logits & $[B,T,V]$ & Vocabulary scores. \\
GPT-2 QKV weight & $[3H,H]$ & Fused projection. \\
GPT-2 MLP input weight & $[F,H]$ & Feed-forward expansion. \\
GPT-2 MLP output weight & $[H,F]$ & Feed-forward contraction. \\
\bottomrule
\end{longtable}}


\begin{intuitionbox}{Intuition}
A shape cheat sheet is a debugging tool.
\end{intuitionbox}

\begin{definitionbox}{Mathematical or contract statement}
Phase-0 GPT-2 has \(H_{kv}=H_q\); later MQA/GQA may not.
\end{definitionbox}

\begin{implementationbox}{Implementation interpretation}
Use the table while reviewing code for transposes.
\end{implementationbox}

\begin{verificationbox}{How to verify this works}
\begin{itemize}
\item Check shapes and dtypes at the boundary.
\item Compare deterministic outputs against PyTorch golden vectors.
\item Record tolerance, seed, model artifact, and backend.
\item Do not benchmark until validation passes.
\end{itemize}
\end{verificationbox}

\begin{warningbox}{Common failure modes}
\begin{itemize}
\item Letting framework defaults choose dtype or layout.
\item Testing only a batch size of one.
\item Depending on upstream checkpoint names in runtime code.
\item Treating benchmark output as validation.
\end{itemize}
\end{warningbox}

\begin{chapterexercises}
\item State the central invariant in one sentence.\\
\textit{Hint.} Use the conceptual guide vocabulary.
\item Construct a toy example with small shapes.\\
\textit{Hint.} Use $B=2$, $T=3$, $H=8$ where possible.
\item Name a validation test that would catch a plausible bug.\\
\textit{Hint.} Prefer generated golden vectors to handwritten long arrays.
\end{chapterexercises}

\chapter{CUDA technical details}
\label{app:cuda-details}
This appendix lists CUDA details that matter for later kernels. \citep{cudaProgrammingGuide}



\begin{intuitionbox}{Intuition}
The bottleneck determines kernel design.
\end{intuitionbox}

\begin{definitionbox}{Mathematical or contract statement}
Coalescing, occupancy, latency hiding, register pressure, shared memory, PTX, SASS, streams, events, and graphs interact.
\end{definitionbox}

\begin{implementationbox}{Implementation interpretation}
Use profiler evidence before architecture-specific tuning.
\end{implementationbox}

\begin{verificationbox}{How to verify this works}
\begin{itemize}
\item Check shapes and dtypes at the boundary.
\item Compare deterministic outputs against PyTorch golden vectors.
\item Record tolerance, seed, model artifact, and backend.
\item Do not benchmark until validation passes.
\end{itemize}
\end{verificationbox}

\begin{warningbox}{Common failure modes}
\begin{itemize}
\item Letting framework defaults choose dtype or layout.
\item Testing only a batch size of one.
\item Depending on upstream checkpoint names in runtime code.
\item Treating benchmark output as validation.
\end{itemize}
\end{warningbox}

\begin{chapterexercises}
\item State the central invariant in one sentence.\\
\textit{Hint.} Use the conceptual guide vocabulary.
\item Construct a toy example with small shapes.\\
\textit{Hint.} Use $B=2$, $T=3$, $H=8$ where possible.
\item Name a validation test that would catch a plausible bug.\\
\textit{Hint.} Prefer generated golden vectors to handwritten long arrays.
\end{chapterexercises}

\chapter{Rust ownership patterns for inference engines}
\label{app:rust-patterns}
Rust patterns include owned buffers, borrowed tensor views, typed errors, and explicit cache mutation. \citep{rustBook}



\begin{intuitionbox}{Intuition}
A view is safe only when storage outlives it.
\end{intuitionbox}

\begin{definitionbox}{Mathematical or contract statement}
Use lifetimes and typed wrappers for shapes, offsets, and byte counts.
\end{definitionbox}

\begin{implementationbox}{Implementation interpretation}
Keep unsafe code out of the phase-0 correctness path.
\end{implementationbox}

\begin{verificationbox}{How to verify this works}
\begin{itemize}
\item Check shapes and dtypes at the boundary.
\item Compare deterministic outputs against PyTorch golden vectors.
\item Record tolerance, seed, model artifact, and backend.
\item Do not benchmark until validation passes.
\end{itemize}
\end{verificationbox}

\begin{warningbox}{Common failure modes}
\begin{itemize}
\item Letting framework defaults choose dtype or layout.
\item Testing only a batch size of one.
\item Depending on upstream checkpoint names in runtime code.
\item Treating benchmark output as validation.
\end{itemize}
\end{warningbox}

\begin{chapterexercises}
\item State the central invariant in one sentence.\\
\textit{Hint.} Use the conceptual guide vocabulary.
\item Construct a toy example with small shapes.\\
\textit{Hint.} Use $B=2$, $T=3$, $H=8$ where possible.
\item Name a validation test that would catch a plausible bug.\\
\textit{Hint.} Prefer generated golden vectors to handwritten long arrays.
\end{chapterexercises}

\chapter{C++ memory and RAII patterns}
\label{app:cpp-patterns}
C++ patterns include RAII file buffers, non-owning spans, and checked tensor views. \citep{cpp2020}



\begin{intuitionbox}{Intuition}
RAII makes lifetime visible.
\end{intuitionbox}

\begin{definitionbox}{Mathematical or contract statement}
An owned buffer can safely produce tensor views if view lifetime is controlled.
\end{definitionbox}

\begin{implementationbox}{Implementation interpretation}
Use sanitiser builds and deterministic tests.
\end{implementationbox}

\begin{verificationbox}{How to verify this works}
\begin{itemize}
\item Check shapes and dtypes at the boundary.
\item Compare deterministic outputs against PyTorch golden vectors.
\item Record tolerance, seed, model artifact, and backend.
\item Do not benchmark until validation passes.
\end{itemize}
\end{verificationbox}

\begin{warningbox}{Common failure modes}
\begin{itemize}
\item Letting framework defaults choose dtype or layout.
\item Testing only a batch size of one.
\item Depending on upstream checkpoint names in runtime code.
\item Treating benchmark output as validation.
\end{itemize}
\end{warningbox}

\begin{chapterexercises}
\item State the central invariant in one sentence.\\
\textit{Hint.} Use the conceptual guide vocabulary.
\item Construct a toy example with small shapes.\\
\textit{Hint.} Use $B=2$, $T=3$, $H=8$ where possible.
\item Name a validation test that would catch a plausible bug.\\
\textit{Hint.} Prefer generated golden vectors to handwritten long arrays.
\end{chapterexercises}

\chapter{PyTorch reference implementation patterns}
\label{app:pytorch-patterns}
PyTorch reference patterns favour clarity, explicit dtype, and stable exports. \citep{pytorchDocs}



\begin{intuitionbox}{Intuition}
Reference code should be easy to distrust productively.
\end{intuitionbox}

\begin{definitionbox}{Mathematical or contract statement}
Use internal tensor names, CPU float32, named intermediates, and deterministic vector output.
\end{definitionbox}

\begin{implementationbox}{Implementation interpretation}
Do not hide contracts in upstream model classes.
\end{implementationbox}

\begin{verificationbox}{How to verify this works}
\begin{itemize}
\item Check shapes and dtypes at the boundary.
\item Compare deterministic outputs against PyTorch golden vectors.
\item Record tolerance, seed, model artifact, and backend.
\item Do not benchmark until validation passes.
\end{itemize}
\end{verificationbox}

\begin{warningbox}{Common failure modes}
\begin{itemize}
\item Letting framework defaults choose dtype or layout.
\item Testing only a batch size of one.
\item Depending on upstream checkpoint names in runtime code.
\item Treating benchmark output as validation.
\end{itemize}
\end{warningbox}

\begin{chapterexercises}
\item State the central invariant in one sentence.\\
\textit{Hint.} Use the conceptual guide vocabulary.
\item Construct a toy example with small shapes.\\
\textit{Hint.} Use $B=2$, $T=3$, $H=8$ where possible.
\item Name a validation test that would catch a plausible bug.\\
\textit{Hint.} Prefer generated golden vectors to handwritten long arrays.
\end{chapterexercises}

\chapter{Artifact validation checklist}
\label{app:artifact-check}
This checklist validates internal artifacts before model execution.


\begin{enumerate}[label=\textbf{G.\arabic*.}]
\item Parse \file{config.json}; reject invalid fields.
\item Parse \file{metadata.json}; check format version, endianness, dtype, and filenames.
\item Parse \file{weights.index.json}; do not rely on key order.
\item Check every required tensor name, dtype, shape, offset, and byte count.
\item Reject byte ranges that overlap or exceed \file{weights.bin}.
\item Verify phase-0 little-endian row-major \dtype{float32}.
\end{enumerate}


\begin{intuitionbox}{Intuition}
Bad artifacts should fail early.
\end{intuitionbox}

\begin{definitionbox}{Mathematical or contract statement}
Validation covers config, metadata, index entries, byte ranges, dtype, endianness, and required names.
\end{definitionbox}

\begin{implementationbox}{Implementation interpretation}
Automate the checklist in every language.
\end{implementationbox}

\begin{verificationbox}{How to verify this works}
\begin{itemize}
\item Check shapes and dtypes at the boundary.
\item Compare deterministic outputs against PyTorch golden vectors.
\item Record tolerance, seed, model artifact, and backend.
\item Do not benchmark until validation passes.
\end{itemize}
\end{verificationbox}

\begin{warningbox}{Common failure modes}
\begin{itemize}
\item Letting framework defaults choose dtype or layout.
\item Testing only a batch size of one.
\item Depending on upstream checkpoint names in runtime code.
\item Treating benchmark output as validation.
\end{itemize}
\end{warningbox}

\begin{chapterexercises}
\item State the central invariant in one sentence.\\
\textit{Hint.} Use the conceptual guide vocabulary.
\item Construct a toy example with small shapes.\\
\textit{Hint.} Use $B=2$, $T=3$, $H=8$ where possible.
\item Name a validation test that would catch a plausible bug.\\
\textit{Hint.} Prefer generated golden vectors to handwritten long arrays.
\end{chapterexercises}

\chapter{Golden-vector generation checklist}
\label{app:golden-check}
This checklist summarises vector generation.


\begin{enumerate}[label=\textbf{H.\arabic*.}]
\item Select authoritative input ids.
\item Run the PyTorch CPU float32 reference.
\item Export JSON sidecar with model identity, dtype, tolerance, and check names.
\item Export NPZ arrays for larger tensors.
\item Store vectors under \file{artifacts/test_vectors/gpt2/gpt2/}.
\item Never paste long arrays into prose or comments.
\end{enumerate}


\begin{intuitionbox}{Intuition}
A reference vector is generated, not invented.
\end{intuitionbox}

\begin{definitionbox}{Mathematical or contract statement}
A vector records input ids, dtype, tolerance, named arrays, and model identity.
\end{definitionbox}

\begin{implementationbox}{Implementation interpretation}
Regenerate vectors only after intentional reference or contract changes.
\end{implementationbox}

\begin{verificationbox}{How to verify this works}
\begin{itemize}
\item Check shapes and dtypes at the boundary.
\item Compare deterministic outputs against PyTorch golden vectors.
\item Record tolerance, seed, model artifact, and backend.
\item Do not benchmark until validation passes.
\end{itemize}
\end{verificationbox}

\begin{warningbox}{Common failure modes}
\begin{itemize}
\item Letting framework defaults choose dtype or layout.
\item Testing only a batch size of one.
\item Depending on upstream checkpoint names in runtime code.
\item Treating benchmark output as validation.
\end{itemize}
\end{warningbox}

\begin{chapterexercises}
\item State the central invariant in one sentence.\\
\textit{Hint.} Use the conceptual guide vocabulary.
\item Construct a toy example with small shapes.\\
\textit{Hint.} Use $B=2$, $T=3$, $H=8$ where possible.
\item Name a validation test that would catch a plausible bug.\\
\textit{Hint.} Prefer generated golden vectors to handwritten long arrays.
\end{chapterexercises}

\chapter{Benchmark checklist}
\label{app:benchmark-check}
This checklist keeps performance measurement separate from correctness.


\begin{enumerate}[label=\textbf{I.\arabic*.}]
\item Run validation first.
\item Record backend, model artifact, prompt length, decode length, and batch size.
\item Warm up before measuring.
\item Separate prefill and decode metrics.
\item Report enough context to reproduce.
\end{enumerate}


\begin{intuitionbox}{Intuition}
A benchmark is meaningful only for a specified workload.
\end{intuitionbox}

\begin{definitionbox}{Mathematical or contract statement}
Benchmark reports should include workload, backend, and validation status.
\end{definitionbox}

\begin{implementationbox}{Implementation interpretation}
Use the same benchmark suite across languages.
\end{implementationbox}

\begin{verificationbox}{How to verify this works}
\begin{itemize}
\item Check shapes and dtypes at the boundary.
\item Compare deterministic outputs against PyTorch golden vectors.
\item Record tolerance, seed, model artifact, and backend.
\item Do not benchmark until validation passes.
\end{itemize}
\end{verificationbox}

\begin{warningbox}{Common failure modes}
\begin{itemize}
\item Letting framework defaults choose dtype or layout.
\item Testing only a batch size of one.
\item Depending on upstream checkpoint names in runtime code.
\item Treating benchmark output as validation.
\end{itemize}
\end{warningbox}

\begin{chapterexercises}
\item State the central invariant in one sentence.\\
\textit{Hint.} Use the conceptual guide vocabulary.
\item Construct a toy example with small shapes.\\
\textit{Hint.} Use $B=2$, $T=3$, $H=8$ where possible.
\item Name a validation test that would catch a plausible bug.\\
\textit{Hint.} Prefer generated golden vectors to handwritten long arrays.
\end{chapterexercises}

\chapter{Glossary}
\label{app:glossary}
The glossary defines recurring terms in this project.


\begin{description}[leftmargin=!,labelwidth=28mm]
\item[Artifact] Internal model directory containing config, metadata, index, and binary weights.
\item[Cache equivalence] Full recomputation equals prefill plus cached decode within tolerance.
\item[Conceptual guide] Repository-level design narrative that explains intent while executable specs define exact contracts.
\item[Decode] Incremental generation path using cached state.
\item[GQA] Grouped-query attention, where \(H_q/H_{kv}\) query heads share each key/value head.
\item[Golden vector] Generated validation record consumed by all runtimes.
\item[KV cache] Per-layer keys and values with logical layout \(\kvshape\).
\item[MQA] Multi-query attention, the \(H_{kv}=1\) special case of grouped key/value sharing.
\item[Prefill] Prompt-processing path that initialises cache state.
\item[RoPE] Rotary positional embedding.
\item[Tensor Core] NVIDIA matrix-multiply acceleration hardware for supported operations.
\end{description}


\begin{intuitionbox}{Intuition}
Precise words prevent imprecise code.
\end{intuitionbox}

\begin{definitionbox}{Mathematical or contract statement}
A term is defined by its contract and verification strategy.
\end{definitionbox}

\begin{implementationbox}{Implementation interpretation}
Extend the glossary when new roadmap phases are added.
\end{implementationbox}

\begin{verificationbox}{How to verify this works}
\begin{itemize}
\item Check shapes and dtypes at the boundary.
\item Compare deterministic outputs against PyTorch golden vectors.
\item Record tolerance, seed, model artifact, and backend.
\item Do not benchmark until validation passes.
\end{itemize}
\end{verificationbox}

\begin{warningbox}{Common failure modes}
\begin{itemize}
\item Letting framework defaults choose dtype or layout.
\item Testing only a batch size of one.
\item Depending on upstream checkpoint names in runtime code.
\item Treating benchmark output as validation.
\end{itemize}
\end{warningbox}

\begin{chapterexercises}
\item State the central invariant in one sentence.\\
\textit{Hint.} Use the conceptual guide vocabulary.
\item Construct a toy example with small shapes.\\
\textit{Hint.} Use $B=2$, $T=3$, $H=8$ where possible.
\item Name a validation test that would catch a plausible bug.\\
\textit{Hint.} Prefer generated golden vectors to handwritten long arrays.
\end{chapterexercises}

\chapter{Bibliography notes}
\label{app:bib-notes}
The bibliography cites foundational papers, model references, and official documentation.



\begin{intuitionbox}{Intuition}
References guide design; validation proves implementation.
\end{intuitionbox}

\begin{definitionbox}{Mathematical or contract statement}
Use papers for architecture ideas and official docs for APIs and version-sensitive details.
\end{definitionbox}

\begin{implementationbox}{Implementation interpretation}
Check official CUDA, PyTorch, Rust, C++, and SafeTensors docs when versions change.
\end{implementationbox}

\begin{verificationbox}{How to verify this works}
\begin{itemize}
\item Check shapes and dtypes at the boundary.
\item Compare deterministic outputs against PyTorch golden vectors.
\item Record tolerance, seed, model artifact, and backend.
\item Do not benchmark until validation passes.
\end{itemize}
\end{verificationbox}

\begin{warningbox}{Common failure modes}
\begin{itemize}
\item Letting framework defaults choose dtype or layout.
\item Testing only a batch size of one.
\item Depending on upstream checkpoint names in runtime code.
\item Treating benchmark output as validation.
\end{itemize}
\end{warningbox}

\begin{chapterexercises}
\item State the central invariant in one sentence.\\
\textit{Hint.} Use the conceptual guide vocabulary.
\item Construct a toy example with small shapes.\\
\textit{Hint.} Use $B=2$, $T=3$, $H=8$ where possible.
\item Name a validation test that would catch a plausible bug.\\
\textit{Hint.} Prefer generated golden vectors to handwritten long arrays.
\end{chapterexercises}
